% ============================================================
%  THE ORTYX PROTOCOL
%  Appendix A: Projection Operators and
%              Constraint-Preserving Reduction
% ============================================================

\chapter{Projection Operators and Constraint-Preserving Reduction}
\label{app:projection}

\noindent
This appendix develops the formal theory of projection
operators and constraint-preserving reduction that underlies
the \Ortyx{} Protocol's treatment of projection dependence
($\AuditP$) and the reinterpretation functor $\RSVPfunctor$.
The formalism is the mathematical backbone of the audit
procedure's ability to distinguish claims that are made
about a full space $X$ from claims that are only established
about a projected representation $M$.

\section{The Basic Setup}
\label{app-a:setup}

Let $X$ denote a high-dimensional trajectory manifold
representing the full dynamical state of a system under
analysis. Let $M$ denote a reduced representational manifold
generated by a projection operator:
\[
  \pi : X \to M.
\]
The projection reduces the dimensionality of the representation,
discarding structure that is not captured in $M$. The
central question of the audit procedure is: which constraints
operative in $X$ survive the projection to $M$?

A representation is said to preserve admissible structure
if there exist constraint operators
\[
  \ConstrX : X \to \mathbb{R}_{\ge 0}, \qquad
  \ConstrM : M \to \mathbb{R}_{\ge 0}
\]
such that for all dynamically admissible trajectories $x \in X_{\mathrm{adm}}$:
\[
  \ConstrM(\pi(x)) \approx \ConstrX(x).
\]
The degree of approximation is measured by the projection
defect functional:
\[
  \ProjDef(x) = \abs{\ConstrX(x) - \ConstrM(\pi(x))}.
\]

\begin{definition}
A reduction $\pi : X \to M$ is \emph{constraint-faithful}
on the admissible trajectory space $X_{\mathrm{adm}} \subseteq X$
if:
\[
  \sup_{x \in X_{\mathrm{adm}}} \ProjDef(x) < \varepsilon
\]
for some specified tolerance $\varepsilon > 0$.
\end{definition}

High $\AuditP$ scores in the \Ortyx{} audit correspond
to frameworks where the projection defect is large relative
to the claims being made: the framework claims to be
establishing properties of $X$ while its evidence is
only about $M$, and $\pi$ is not demonstrated to be
constraint-faithful with respect to those properties.

\section{The Projection Trap}
\label{app-a:projection-trap}

A particular failure mode occurs when optimization
acts directly on $M$ while feedback from $X$ becomes
inaccessible. This is the \emph{projection trap}:
\[
  \frac{d}{dt}\,\mathrm{Dist}\!\bigl(x(t),\,
  \pi^{-1}(m(t))\bigr) \to \infty.
\]
The trajectory in $X$ and the trajectory in $M$ drift apart:
the representational dynamics diverge from the full-space
dynamics. This produces:

\begin{itemize}
  \item \textit{Proxy divergence}: optimization on $M$
  fails to track the actual objective in $X$.
  \item \textit{Representational collapse}: the map
  $\pi^{-1}(m(t))$ --- the preimage of the current
  representation --- shrinks to an increasingly unrepresentative
  subset of $X_{\mathrm{adm}}$.
  \item \textit{Semantic drift}: vocabulary calibrated
  to $M$ becomes increasingly disconnected from the
  phenomena in $X$ it was originally designed to describe.
\end{itemize}

In the context of the \Ortyx{} audit, semantic drift is
the mechanism by which vocabulary migration (identified
in Chapter~6) produces $\FailGeo{3}$ (Semantic Universality
Collapse): the vocabulary, optimized against its usage
in $M$, acquires meanings incompatible with its intended
referents in $X$.

\section{Multi-Level Projections and the Projection Ladder}
\label{app-a:projection-ladder}

In the RSVP framework, the relationship between levels
of description is formalized as a projection ladder:
a sequence of projections
\[
  X = X_0 \xrightarrow{\pi_1} X_1
  \xrightarrow{\pi_2} X_2 \xrightarrow{\pi_3}
  \cdots \xrightarrow{\pi_k} X_k = M
\]
where each $\pi_i$ is a constraint-faithful reduction
at a specified tolerance $\varepsilon_i$. The accumulated
projection defect of the full ladder is:
\[
  \Delta_{\pi,\,\mathrm{total}} \le \sum_{i=1}^{k} \varepsilon_i
\]
by the triangle inequality, assuming the defects compose
in the worst case. Constraint-faithful projections at
each rung of the ladder bound the total information
loss from the full space to the coarsest representation.

\begin{definition}
A projection ladder $(\pi_1, \ldots, \pi_k)$ is
\emph{cumulatively constraint-faithful} at tolerance
$\varepsilon$ if $\sum_{i=1}^k \varepsilon_i < \varepsilon$.
\end{definition}

The RSVP reinterpretation functor $\RSVPfunctor$
applied in Chapter~16 is, in this formalism, a
rung of the projection ladder: it maps Phoenix concepts
to RSVP-compatible concepts in a way that is designed
to be constraint-faithful at the level of the engineering
results while acknowledging the accumulated defect
from engineering to philosophical levels.

\section{Constraint-Faithfulness Tests}
\label{app-a:faithfulness-tests}

For a projection $\pi : X \to M$ to be demonstrably
constraint-faithful, the following must be established:

\begin{enumerate}
  \item \textit{Specification of $\ConstrX$ and $\ConstrM$}:
  Both constraint operators must be operationally defined
  --- computable from observations in $X$ and $M$
  respectively.
  \item \textit{Bounding of $\ProjDef$}: An upper bound
  on $\ProjDef(x)$ across $X_{\mathrm{adm}}$ must be
  derived from the properties of $\pi$ and the constraints.
  \item \textit{Admissibility characterization}: The
  admissible trajectory space $X_{\mathrm{adm}}$ must
  be characterized sufficiently to determine whether
  the bound on $\ProjDef$ holds throughout it.
\end{enumerate}

The $\AuditP$ operator of the \Ortyx{} Protocol assigns
high scores to frameworks that fail any of these three
requirements for the specific projections their claims
depend on.
